\scanpage{152}
\mspage{77}
% archive sheet number: 151
% L1–2
\diagram{\[
\begin{tikzcd}[ampersand replacement=\&]
G \arrow[r,"\chi"] \arrow[dr,"\chi'"'] \& (\mathbb Z/n\mathbb Z)^* \arrow[d,"\mathrm{can.}"]\\
\& Z\simeq(\mathbb Z/n\mathbb Z)^*/\pm1\ .
\end{tikzcd}\tag*{(59)}
\]}
% L3–5
On voit que le noyau de $\chi$ est $G_0^+$, de sorte que la suite exacte (17) peut aussi se déduire d'une suite exacte
% L6–7
\diagram{\[
\begin{tikzcd}[ampersand replacement=\&,column sep=1.3em]
1 \arrow[r] \& G_0^+ \arrow[r,dash] \& G \arrow[r,"\chi","\text{car. cyclotomique}"'] \& (\mathbb Z/n\mathbb Z)^* \arrow[r] \& 1\ .\\
\& \& Z\simeq(\mathbb Z/n\mathbb Z)^*/\pm1 \arrow[u,hook] \arrow[ur,maps to,"\dot\lambda\mapsto\lambda^2"'] \& \&
\end{tikzcd}\tag*{(60)}
\]}
% slanted
\margin{sauf dans cas octaédral ($n=4$) où il faut remplacer $Z$ par $\{1\}$}
% L8–10
\deleted{\uline{NB} On lit} par
\[
\begin{aligned}
G_0&=\{g\in G\mid\chi(g)\in\{\pm1\}\}\\
&[=\operatorname{Ker}(\chi':G\longrightarrow(\mathbb Z/n\mathbb Z)^*/\pm1\ ]
\end{aligned}\tag*{(61)}
\]

% L11–12
\uline{NB} On lit sur (60) que
\[
\underbrace{Z\cap G_0^+}_{=\mathrm{Cent}(G_0^+)}=
({}_2(\mathbb Z/n\mathbb Z)^*)/(\pm1)\simeq\mathbb F_2^{\mathbb P(n)'}/\mathrm{diag}
\tag*{(62)}
\]
% slanted
\margin{y compris cas $n\deleted{\unsure{\ne}}=4$, \deleted{Dans le cas $n=4$} où \deleted{\unsure{$Z$ est de}} on trouve \deleted{\unsure{où}} $\underbrace{Z}_{\{1,\unsure{a|0}\}}\cap G_0^+=\{1\}$}
% L13–16
$\mathbb P(n)$ est l'ensemble des nbs premiers \deleted{impairs} qui divisent $n$, et $\mathbb P'(n)$ est égal à $\mathbb P^{\deleted{\unsure{*}}}(n)$ sauf si $2\in\mathbb P(n)$ et si 2 intervient avec un \deleted{\unsure{fact}} exposant \inserted{$e_2$} $\ne2$ \deleted{i.e. sauf si $\left\{\begin{smallmatrix}n\equiv0\ (2)\\ n\not\equiv0\ (4)\end{smallmatrix}\right.$, auquel cas} auquel cas on pose
% L17–19
\[
\left\{\begin{aligned}
\mathbb P'(n)&=\mathbb P(n)\setminus\{2\}&&\text{si }e_2=1\\
\mathbb P'(n)&=\mathbb P(n)\amalg\{\substack{\text{un élément}\\\text{supplém.}}\}&&\text{si }e_2\ge3\\
(\mathbb P'(n)&=\mathbb P(n)&&\text{si }e_2\in\{0,2\}
\end{aligned}\right.\tag*{(63)}
\]
% L20–21
donc
\[
Z\cap G_0^+\simeq\mathbb F_2^\nu\quad\text{avec }\nu=
\underbrace{\nu(n)}_{\substack{\text{nb de facteurs}\\\text{premiers de }n}}
\unsure{-|+}1+\varepsilon,\quad\varepsilon=
\begin{cases}
\unsure{0}\ \deleted{\unsure{\text{si}}}\ e_2\in\{0,\unsure{2|1}\}\\
-1\ \text{si }e_2=1\\
+1\ \text{si }e_2\ge3
\end{cases}.
\tag*{(64)}
\]
